\section*{Hoofdstuk \ref{ch:g5:irreversible-changes} --- Veranderingen die je niet kunt terugdraaien}

\begin{solution}{exo:g5:irreversible-changes:1}
Chocolade smelten: terug te draaien (ze wordt weer hard als ze afkoelt).
Een lucifer \omterm{def:g4:what-a-fire-needs:burning}{verbranden}: niet. Zout \omterm{def:g2:mixing-and-dissolving:dissolve}{oplossen}: terug te draaien (laat het
water verdampen). Brood roosteren: niet.
\end{solution}

\begin{solution}{exo:g5:irreversible-changes:2}
Een verandering waarbij nieuwe stoffen ontstaan. In de keuken: een ei
bakken, een cake bakken (ook: brood roosteren, een doorgesneden appel die
bruin wordt).
\end{solution}

\begin{solution}{exo:g5:irreversible-changes:3}
Door het water te laten verdampen: het zout blijft in het schaaltje.
\omterm{def:g2:mixing-and-dissolving:dissolve}{Oplossen} is een \omterm{def:g5:irreversible-changes:reversible}{omkeerbare verandering}.
\end{solution}

\begin{solution}{exo:g5:irreversible-changes:4}
Een nieuwe geur, en klontjes (een vaste stof) die in de melk verschijnen.
\end{solution}

\begin{solution}{exo:g5:irreversible-changes:5}
De veranderingen links: suiker die in water \omterm{def:g2:mixing-and-dissolving:dissolve}{oplost}, ijs dat smelt, zand
en grind die gemengd worden. De dubbele pijl betekent dat de verandering
twee kanten op kan: je kunt ze terugdraaien.
\end{solution}

\begin{solution}{exo:g5:irreversible-changes:6}
De olie houdt \omterm{def:g4:air-a-mixture-of-gases:air}{lucht} en water weg van het ijzer van de ketting: ze
vertraagt het roesten.
\end{solution}

\begin{solution}{exo:g5:irreversible-changes:7}
Gasbelletjes die in een vloeistof verschijnen zonder dat er verwarmd
wordt.
\end{solution}

\begin{solution}{exo:g5:irreversible-changes:8}
Nuttig: eten koken, brood bakken, gas \omterm{def:g4:what-a-fire-needs:burning}{verbranden} om te verwarmen.
Schadelijk: roesten, eten dat bederft, hout dat rot.
\end{solution}

\begin{solution}{exo:g5:irreversible-changes:9}
De kou vertraagt de \omterm{def:g5:irreversible-changes:chemical-change}{chemische veranderingen} waardoor melk zuur wordt,
dus blijft ze langer goed.
\end{solution}

\begin{solution}{exo:g5:irreversible-changes:10}
Nee. De belletjes zijn waterdamp, die weer water wordt als hij afkoelt:
er ontstaat niets nieuws, en de verandering kun je terugdraaien.
Belletjes alleen bewijzen niet dat er een nieuwe stof is ontstaan.
\end{solution}

\begin{solution}{exo:g5:irreversible-changes:11}
Een brandende kaars geeft warmte en licht af, en maakt nieuwe stoffen:
waterdamp (een koud glas erboven beslaat) en koolstofdioxide. Het
kaarsvet dat verbrandt, komt niet terug. Gesmolten kaarsvet daarentegen
wordt weer hard als het afkoelt.
\end{solution}
